\section{Boundary incidence does not resolve all positive Q conflicts}
\label{sec:boundary-transversal}

The parameter $h(T)$ of Section~\ref{sec:euler-transversal} depends on
the complete positive root-ray conflict hypergraph, not merely on the
boundary size.  A finite one-vertex torus triangulation has only two
boundary triangles and three boundary edges.  This suggests two potential
shortcuts: resolve the tetrahedra containing boundary faces, or resolve
all tetrahedra incident to boundary edges.  Both coverage statements are
false.  This section supplies exact counterexamples, including one with
eleven tetrahedra, together with a solver-free replay from original normal
matching equations.  The counterexamples establish neither a family with
growing $h(T)$ nor a lower bound on general recognition complexity.

\subsection{The two proposed coverage statements}

Write $K(T)$ for the full root Q cone, $F$ for canonical formal Euler,
and $D(q)$ for the complete set of tetrahedra supporting at least two
positive Q types.  Define
\[
 B_{\mathrm{face}}(T)=\{\tau:\tau\text{ contains a boundary face}\},
 \qquad
 B_{\mathrm{edge}}(T)=\{\tau:\tau\text{ contains a boundary edge}\}.
\]
Here ``contains'' refers to an incidence in the generalized triangulation;
the same global edge may occur several times in a tetrahedron.  Certainly
$B_{\mathrm{face}}\subseteq B_{\mathrm{edge}}$.

For a one-vertex torus boundary, $3f=2e$ and $1-e+f=0$, giving
$f=2$ and $e=3$.  Thus $|B_{\mathrm{face}}|\le2$.  This counting
argument places no constant bound on $|B_{\mathrm{edge}}|$, since a
boundary edge may be incident to many tetrahedra.  More fundamentally,
neither set necessarily intersects every $D(r)$ for incompatible root
extreme rays $r$ with $F(r)>0$.

\begin{proposition}[Failure of boundary conflict coverage]
\label{prop:boundary-conflict-failure}
There are finite, connected, orientable, one-vertex triangulations with
one torus boundary for which $B_{\mathrm{face}}$ is not a transversal
of $\mathcal H(T)$.  The same is true of $B_{\mathrm{edge}}$.
The retained examples arise from knot-exterior sources and explicit
boundary layerings.  No assumption concerning the absence of admissible
canonical positive vectors is needed for these failures.
\end{proposition}

The proof consists of the exact witnesses below.  In particular, their
extremality is checked independently; positivity of a nonextreme formal
point alone would not suffice for this proposition on an arbitrary source.

\subsection{Root-ray extremality from original STANDARD matching}

Let $M$ be the original STANDARD matching matrix.  For a set $S$ of Q
coordinates, let $M_S$ retain all $4t$ triangle columns and only the Q
columns in $S$.  Let $M_\triangle$ retain just the triangle columns.
These matrices are formed directly from face incidences, without the
Q-cycle elimination used by the search oracle.

\begin{lemma}[A supported Q rank certificate]
\label{lem:boundary-original-standard-rank}
For a finite one-vertex source, suppose a nonzero nonnegative Q vector
$q$ with support $S$ has a STANDARD matching lift.  If
\[
 (4t+|S|-\operatorname{rank}M_S)
       -(4t-\operatorname{rank}M_\triangle)=1,
\]
then $q$ spans an extreme ray of the full root Q cone.  In this setting
$4t-\operatorname{rank}M_\triangle=1$.
\end{lemma}
\begin{proof}
With Q coordinates zero, each face equation equates triangle values at
the two corresponding corners.  Their equivalence classes are exactly
the source vertices.  Consequently the triangle-only matching kernel
is the one-dimensional vertex-link line.

Project $\ker M_S$ onto its Q coordinates.  Its kernel is precisely
this triangle-only line, so the displayed difference is the dimension
of the projected supported space.  Any real matching lift can be made
nonnegative in its triangle coordinates by adding a sufficiently large
multiple of the vertex-link vector; this does not change its Q vector.
Thus projection and restriction describe the supported face of the
nonnegative root Q cone.  The given nonzero nonnegative vector lies in
that face.  If the projected space has dimension one, the face is a
single ray.  Since vanishing coordinate constraints define a face of
a nonnegative cone, that ray is also extreme in the root cone.
\end{proof}

The replay checks $\operatorname{rank}M_\triangle$ explicitly in
addition to using the source's one-vertex validation.  All ranks use
exact rational elimination.  It also checks every original matching
equation, nonnegative integer coordinates, Q-coordinate gcd one, and
zero minimum triangle coordinate over the unique vertex class.
The last condition identifies the canonical lift.

Formal Euler is computed directly as the number of intersection vertices
minus normal arcs plus polygon pieces.  For a local edge $ab$, its
intersection weight is
\[
 x_a+x_b+\sum_{i=0}^2q_i-q_{\operatorname{quad}(a,b)}.
\]
The replay checks agreement among all occurrences of each global edge,
counts each global face once, and subtracts its arc count.  Incompatible
coordinates are used only as formal matching vectors; they do not
describe embedded normal surfaces.

\subsection{Boundary-face counterexamples at retained endpoints}

Two labelled endpoint sources are copied verbatim from the previous
transversal audit.  They have the following retained witness, indexed
zero in each fixture.  The listed support rank is the dimension rank in
Q coordinates, derived here through Lemma~\ref{lem:boundary-original-standard-rank}.
\begin{center}
\begin{tabular}{lrrrrl}
source & $t$ & $F$ & $|S|$ & rank & $D(q)$ \\
\hline
trefoil endpoint & 10 & 1 & 10 & 9 & $\{0,5\}$ \\
figure-eight endpoint & 13 & 2 & 11 & 10 & $\{3,8,12\}$
\end{tabular}
\end{center}
The trefoil boundary faces are $(6,0),(9,0)$, so its conflict set misses
$B_{\mathrm{face}}=\{6,9\}$.  The figure-eight faces are $(7,0),(9,0)$,
and its conflict set misses $\{7,9\}$.  In fact, all eight retained
trefoil rays and nine of the eleven retained figure-eight rays miss
their respective boundary-face sets.  These are finite checks of
retained discoveries, not an exhaustive enumeration of root rays.

The corresponding canonical JSON source SHA256 digests are
\begin{quote}\small\ttfamily
trefoil:\newline
515b94f650b5a7b0e8e09322ba31f59eb8cdb108fe33d9220994c43f148ca3fd\newline
figure-eight:\newline
1d03c1400d5168f3f862ddbccaa7551650a95f12d850d563fd9de653f0864c38
\end{quote}
These labelled equalities matter: a knot name or tetrahedron count would
not identify the same conflict hypergraph.  At these original endpoints,
$B_{\mathrm{edge}}$ contains all ten trefoil tetrahedra and all
figure-eight tetrahedra except $1$ and $5$.  The original observations
alone therefore do not establish the edge-star counterexample.

\subsection{A single boundary layer gives the edge-star refutation}

Start from the ten-tetrahedron trefoil source above.  Add tetrahedron
$10$, replacing its two old boundary records by
\[
 (6,0)\longrightarrow(10,[3,1,2,0]),\qquad
 (9,0)\longrightarrow(10,[2,3,0,1]).
\]
The bracket lists the image of local vertices $0,1,2,3$.
New faces $0$ and $1$ are boundary; new face $2$ is paired to tetrahedron
$9$ with permutation $[2,3,0,1]$, and new face $3$ to tetrahedron $6$
with permutation $[3,1,2,0]$.  All other old face pairings remain unchanged.
The edge being layered on is local edge $4$ of tetrahedron $6$, where
the local ordering is $01,02,03,12,13,23$.

The primitive witness's STANDARD rows, followed by the new row, are
\[
\begin{array}{c|rrrr|rrr}
\tau&T_0&T_1&T_2&T_3&Q_0&Q_1&Q_2\\\hline
0&4&0&2&8&0&2&2\\
1&1&4&3&2&0&5&0\\
2&6&4&2&10&0&0&0\\
3&3&1&10&2&3&0&0\\
4&3&1&2&8&3&0&0\\
5&3&0&2&1&9&1&0\\
6&6&2&10&2&0&0&0\\
7&1&4&3&10&0&3&0\\
8&3&1&2&2&11&0&0\\
9&6&2&2&2&8&0&0\\\hline
10&2&2&10&10&0&0&0
\end{array}
\]
The new source has digest
\begin{quote}\small\ttfamily
f6a05037099e93211b93a865e6ebca666d033e91dc04d9cb17475544206d5daa.
\end{quote}
It is finite, connected, orientable, one-vertex, and has one torus boundary.
The full boundary-edge incidence union is now
\[
 B_{\mathrm{edge}}=\{2,3,6,9,10\},
 \qquad D(q)=\{0,5\},\qquad F(q)=1.
\]
The two sets are disjoint.  There are $44+10=54$ retained STANDARD
columns and their matching rank is $52$.  The triangle block has rank
$43$, giving projected support dimension $54-52-(44-43)=1$.
This proves the second assertion of
Proposition~\ref{prop:boundary-conflict-failure} using a genuine positive
root extreme ray.

\subsection{Why the layer preserves the compact manifold}

A validity flag alone is not a homeomorphism proof.  The relevant local
statement is Jaco--Rubinstein, \emph{Layered-triangulations of 3-manifolds},
Section~3, printed page~9, together with the boundary-flip
discussion in Lemma~9.2, printed pages~85--86:
\url{https://arxiv.org/pdf/math/0603601v2}.
Their generalized triangulations allow identified simplex boundaries.
The local definition attaches two tetrahedron faces to two distinct
boundary triangles consistently along their shared edge.  It explicitly
allows identifications among the sides of the boundary quadrilateral and
states that the resulting triangulation represents the same compact
manifold.  Thus one must not replace this statement by an argument
requiring an embedded attaching disk.

For every retained step, the replay checks the two distinct old boundary
triangles, the reciprocal face permutations, and agreement of the shared
attaching edge including its oriented ends.  It checks that every other
old pairing remains unchanged, no distinct old edge classes merge, and
exactly the chosen old boundary edge is replaced by one new boundary
edge.  Independent native source validation checks manifold links,
orientability, connectedness and torus boundary.  These combinatorial
checks identify the explicit attachment with the cited layering operation;
the cited local theorem supplies the collar homeomorphism inference.

The new layerings are not operations in the maintained crushing-chain
certificate schema.  Accordingly, these are supplied-source
counterexamples with checked collar maps, not a claim that the current
diagram driver emitted the layered sources as authenticated successors.
The initial sources agree exactly with the previously retained
diagram-bound endpoints.  This package does not repeat their discovery
or substitute knot type for a chain certificate.

\subsection{Longer collars and the remaining structural question}

The retained construction also flips each of the three original boundary
edges once.  Trefoil uses original local representatives
$(6,4),(9,4),(6,3)$; figure-eight uses $(7,4),(7,5),(7,3)$.
Every step retains the old coordinates, adds at most one Q type in its
new tetrahedron, and preserves canonical formal Euler and the complete
conflict set.  The resulting exact checks are
\begin{center}
\begin{tabular}{lrrrll}
source & $t$ & $F$ & $|S|$ & $D(q)$ & $B_{\mathrm{edge}}$\\\hline
trefoil, three layers &13&1&12&$\{0,5\}$&$\{10,11,12\}$\\
figure-eight, three layers &16&2&13&$\{3,8,12\}$&$\{13,14,15\}$
\end{tabular}
\end{center}
Their original STANDARD restricted ranks over the three steps are
$52,57,62$ and $65,70,75$, respectively.  Each projected Q support
still has dimension one.  Complete maps and coordinates, including all
intermediate sources, are retained rather than reconstructed from knot
names.

The Luo--Tillmann standard-solution basis is a linear basis with signed
coefficients, not a positivity decomposition that forces conflicts near
the boundary; see Theorem~6, Lemma~7 and the following Euler computation
in \url{https://arxiv.org/pdf/math/0510537v2}.  The examples retain
positive formal Euler and core conflicts while moving all boundary-edge
incidences into a collar.  Therefore that basis cannot imply either
refuted coverage statement under only the stated finite one-vertex
torus-boundary hypotheses.

Even the size of an edge star can grow with collars.  In successive
boundary flips, keep one boundary edge and flip another.  The kept edge
remains in the boundary and is incident to each newly layered tetrahedron.
Its star then grows linearly with the number of layers.  This observation
concerns incidence size, not the minimum resolving number $h(T)$;
simplification could remove the collars.  A useful boundary-based search
theorem would therefore need explicit extra canonicalization or
collar-removal hypotheses and a new coverage argument.  Neither a
uniform small-$h$ theorem nor a growing-$h$ counterfamily is proved here.

\subsection{Retained inputs and fresh solver-free replay}

The package \path{fast/boundary_transversal_research} contains the exact
original endpoint fixture bytes, all layered sources and coordinates,
and the original discovery scripts and outputs with SHA256 provenance.
Historical generation pins are retained as historical evidence; they
are not asserted to equal current source hashes.  The current replay
pins all native Python files and all package sources and immutable data
at both ends of its run.  It checks all nineteen endpoint witnesses and
all six layerings, including first-step edge-star disjointness in the
eleven-tetrahedron example.

Run from \path{fast}:
\begin{verbatim}
python -B -S -m boundary_transversal_research.replay \
  --output /tmp/boundary-transversal-audit.json
\end{verbatim}
The checker uses neither Regina nor LP optimization nor Q-cycle
elimination.  A separate optional \path{generate} module performs fresh
Regina layerings and the small local extension searches, comparing
every labelled result with the retained fixture.  Its bounded fixture
enumeration is not claimed to be a polynomial construction for arbitrary
input coordinate magnitudes.  No timings are reported here.  In
particular, neither the old retained discovery nor fresh replay closes
the general quasi-polynomial recognition gap.

The published \path{synthesis/data/boundary-transversal-audit.json}
records all nineteen ray checks and six layer checks.  The corresponding
\path{boundary-transversal-regeneration.json} records a fresh Regina~7.4
construction agreeing exactly with every retained source and coordinate
vector.  An exact saved-evidence replay and seven passing rejection
controls are retained alongside them.  Audit, regeneration and replay
agree on the same $306$ native/package file hashes; these are correctness
records, with no running-time comparison.
